\documentclass[twocolumn]{aastex631}

\usepackage{amsmath}
\usepackage{multirow}
\usepackage{natbib}
\usepackage{graphicx} 
\usepackage{aas_macros}

\begin{document}

The search for a consistent theory of quantum gravity has led to extensive investigations into modifications of General Relativity, particularly scalar-tensor theories. Among the most promising are the Degenerate Higher-Order Scalar-Tensor (DHOST) theories, which represent the most general class of scalar-tensor models that, despite containing higher-order derivatives, propagate only the three desired degrees of freedom: the two polarizations of the graviton and a single scalar mode. This remarkable feat is accomplished by imposing a set of `degeneracy conditions' on the arbitrary functions defining the Lagrangian. These conditions enforce a delicate cancellation mechanism that eliminates the Ostrogradsky ghost, a pathological instability that generically plagues higher-derivative theories and signals a loss of unitarity. While this classical stability makes DHOST theories compelling frameworks for cosmology and gravitational physics, their ultimate viability rests on whether this crucial degeneracy is preserved under quantum corrections. The central problem, which remains largely open, is whether the intricate algebraic structure that ensures classical health is robust enough to survive the introduction of quantum loops, or if radiative corrections inevitably reintroduce the very instability that was so meticulously removed.Addressing this question for the entirety of the DHOST landscape is a formidable challenge due to the vast parameter space of free functions and the notorious complexity of quantum field theory calculations in curved spacetime. A direct, brute-force analysis is intractable. Therefore, a more strategic approach is required to probe the quantum fate of classical degeneracy. In this paper, we adopt such a strategy by using symmetry as a guiding principle to isolate a controlled and physically motivated theoretical laboratory. We focus on a subclass of Type Ia DHOST theories and investigate the conditions under which they can possess a field-dependent gauge symmetry, characterized by the transformations $\delta_\epsilon \phi = \epsilon \Lambda(\phi, X)$ and a corresponding field-dependent diffeomorphism of the metric. Rather than postulating such a symmetry ad-hoc, we first derive the mathematical integrability conditions that the theory's free functions must satisfy for this symmetry to exist. This procedure rigorously defines a non-trivial, self-consistent subclass of symmetric models, from which we select the simplest representative case for our quantum analysis.Our method is to compute the one-loop effective action for this symmetric model using the path integral formalism and the background field method \citep{ferrero2025oneloopeffectiveactioncoherent,ferrero2025oneloopeffectiveactioncoherent}. By expanding the action to second order in quantum fluctuations of the scalar field and the metric around a classical background, we isolate the quadratic operator governing their dynamics \citep{ferrero2025oneloopeffectiveactioncoherent}. The ultraviolet divergences arising from the functional determinant of this operator are then calculated using the heat kernel expansion \citep{branchina2024pathintegralmeasurecosmological,ferrero2025oneloopeffectiveactioncoherent}. These divergences manifest as local counter-terms that must be added to the classical action \citep{ferrero2025oneloopeffectiveactioncoherent}. The first part of our analysis is to test whether these quantum-generated counter-terms respect the original symmetry. We demonstrate that they do not; the classical symmetry is broken at the quantum level, giving rise to a quantum anomaly. The effective action, corrected by one-loop effects, is no longer invariant under the transformations that defined the classical theory.The central result of this work, however, goes beyond simply identifying an anomaly. We establish a direct and inexorable link between this symmetry breaking and a catastrophic failure of the theory's stability. By scrutinizing the structure of the anomaly-inducing counter-terms, including higher-curvature terms such as a Weyl-squared invariant, we show that they explicitly violate the original DHOST degeneracy conditions. This violation signifies that the delicate cancellation mechanism at the heart of the classical theory has been destroyed by quantum effects. Consequently, the Ostrogradsky ghost degree of freedom is reintroduced into the particle spectrum, rendering the theory quantum-mechanically unstable. This finding constitutes a no-go theorem for this class of symmetric theories: the presence of a classical gauge symmetry that becomes anomalous at the quantum level leads directly to the loss of quantum stability. This elevates the requirement of anomaly cancellation from a mere aesthetic preference to a powerful and necessary selection principle for constructing consistent and viable theories of modified gravity.

\end{document}
                